\scanpage{3}
\mspage{1}
\uline{Théorème} $S$ schéma loc. noeth, $U$ ouvert, $T$ fermé complémentaire, $p$ \unsure{nb|est} premier, tel que $\forall t\in T$, $\mathrm{car}\ k(t)=p$. Considérons la catégorie $C$ des schémas abéliens \inserted{$A$} sur $S$, et la catégorie $C'$ des couples $(A_U,M,\varphi)$, où $A_U$ est un schéma abélien sur $U$, $M$ un \unsure{groupe|gpe} $p$-divisible sur $S$, et $\varphi:M|U\longrightarrow T_p(A_U)$. Soit
\[
F:C\longrightarrow C'
\]
le foncteur
\[
F(A)=(A\times_S U,T_p(A),\mathrm{id}_{T_p(A)|U}).
\]

Alors :

1) Supposons que $U$ rencontre toute composante connexe de $S$ \inserted{(p.ex. $U$ dense dans $S$, i.e. $\operatorname{prof\,ét}_T(S)\geq 1$).} Alors le foncteur $F$ est pleinement fidèle.

2) Supposons que \deleted{$U$ soit dense dans $S$ et que} a) \deleted{\unsure{le morphisme} \illegible{2} \unsure{de} $S$ \unsure{sur} $S$, \unsure{et}} \inserted{$\exists$ morphisme fini surjectif. $S'\longrightarrow S$ avec $S'$ géom. unibranche (p.ex. la normalisée de $S$ soit finie sur $S$, \unsure{où|ou|si} $S$ local).}

\deleted{b) pour tout $t\in T$, si $S'$ est le localisé strict de $S$ en $t$, \inserted{$\bar t$ son point fermé,} et $S'_1$ et $S'_2$ deux composantes irréductibles de $S'$, on \deleted{\unsure{ait}} ait $S'_1\cap S'_2\neq\{\bar t\}$. \deleted{Sous ces conditions, $F$ est une équivalence de catégories.} (N.B. Cette condition équivaut aussi à celle-ci : l'image inverse $U''$ de $U$ dans $S''=S'\times_S S'$ est dense dans $S''$. Elle est vérifiée si $S$ est \deleted{\unsure{géom.}} unibranche en les points de $T$) Alors \unsure{a)} b) Posant} b) Que \unsure{-|.} $\operatorname{prof\,ét}_T(S)\geq 2$\tikz[baseline=0pt,x=1ex,y=1ex]{\draw (0,2.1)--(.06,.20) .. controls (-.03,-.12) and (-.43,-.09) .. (-.28,.16) .. controls (-.2,.31) and (-.04,.27) .. (.06,.20);}. Alors $F$ est \deleted{pleinement fidèle} \inserted{une équivalence de} catégories

\margin{Pour 1) TSVP : $\longrightarrow$}

\tikz[baseline=0pt,x=1ex,y=1ex]{\draw (0,2.1)--(.06,.20) .. controls (-.03,-.12) and (-.43,-.09) .. (-.28,.16) .. controls (-.2,.31) and (-.04,.27) .. (.06,.20);} Soit, avec les notations \deleted{de} \inserted{du} N° a), $S''=S'\times_S S'$, $S'''=S'\times_S S'\times_S S'$. Considérons l'ensemble simplicial $\overset{\substack{L_0\\\text{''}}}{\pi_0(S')}\leftleftarrows\overset{\substack{L_1\\\text{''}}}{\pi_0(S'')}\mathrel{\substack{\leftarrow\\[-.65ex]\leftarrow\\[-.65ex]\leftarrow}}\overset{\substack{L_2\\\text{''}}}{\pi_0(S''')}$, \deleted{\unsure{les|Com}} et \deleted{\unsure{soit}} soit $K_0\leftleftarrows K_1\mathrel{\substack{\leftarrow\\[-.65ex]\leftarrow\\[-.65ex]\leftarrow}} K_2$ le ss-ensemble simplicial formé avec les composantes connexes dont l'image dans $S$ rencontre $U$. Supposons \unsure{:} que : \deleted{pour tout} \inserted{b')} la plus petite partie de $L_1$ contenant $K_1$ et stable par ``troisième face'' soit $L_1$ (c'est ce qui est assuré par la condition b)). Alors $F$ est \deleted{\unsure{pl. fid.|pl. f.}} une équivalence de catégories .....

% archive sheet number: 2
